\documentclass[10pt,a4paper]{amsart}
\usepackage[margin=19mm]{geometry}
\usepackage{amsmath,amssymb,mathtools}
\usepackage[scr=rsfs,cal=euler]{mathalfa}
\usepackage{xfrac}
\usepackage[unicode,hidelinks,bookmarks=false]{hyperref}
\newcommand{\ie}{i.e.\ }
\newcommand{\cf}{cf.\ }
\newcommand{\resp}{respectively\ }
\newcommand{\Subsection}{\subsection}
\providecommand{\nobreakdash}{\nobreak\hbox{-}}
\setlength{\emergencystretch}{2em}
\multlinegap0pt
\newtheorem{proofof}{Proofof}
\usepackage{tikz}
\usepackage[shortlabels]{enumitem}
\usepackage{amssymb}
\usepackage{dsfont}
\usepackage[normalem]{ulem}
\usepackage{xcolor}
\usepackage{mathtools}
\newtheorem{lemma}{Lemma}
\newcommand{\hit}{\tau_{\rm hit}}
\title{Mixing times and spectra of non-equilibrium symmetric exclusion processes on general graphs}
\author{Leonard J. Schulman, Alistair Sinclair}
\date{Source: arXiv:2607.22991v1 (CC BY 4.0)}
\begin{document}
\maketitle

\noindent\emph{Source and licence.} Mixing times and spectra of non-equilibrium symmetric exclusion processes on general graphs. Version arXiv:2607.22991v1 (CC BY 4.0), distributed by arXiv under the Creative Commons Attribution 4.0 International licence, http://creativecommons.org/licenses/by/4.0/, as recorded in the arXiv licence metadata for this exact version and shown on the abstract page https://arxiv.org/abs/2607.22991v1. Full licence notice: \texttt{licenses/arxiv-2607.22991-CC-BY-4.0.txt}.\par\smallskip

\noindent\textit{Editorial scope.} Counted unit: the article's lemma lem:lowerbd (source0.tex lines 528--532). The article's complete original proof of it is delivered with it (source0.tex lines 564--624, one \texttt{proof} environment of 61 lines, which the article places away from the statement). No mathematics is written, repaired or re-derived: the statement and the proof are the article's own lines, in the article's own notation, and no retained line is substituted or re-flowed.  No further source-local context is needed: every cross-reference of every retained line resolves inside the two delivered spans.  Nothing was omitted from any retained span, so the package carries no omission record.  No retained line contains a citation, so the wrapper carries no bibliography and no bibliography entry.  No retained line contains a figure environment, an image-inclusion command or drawing code, so no image is delivered and the package declares no asset dependency.  Editorial formatting only, in addition: an AMS \texttt{amsart} preamble that restates the article's own declarations and macros used by the retained lines, namely the environments \texttt{lemma} on separate counters, together with the standard packages those lines need.  The retained environments print in this document as Lemma 1 in order of appearance, whereas the article prints the same items with the numbering of its own full document.  The complete native source of the article as distributed in the arXiv e-print for this exact version is bundled unchanged as \texttt{sources/205991/source0.tex}.
\begin{lemma}\label{lem:lowerbd}
There exists a subset $S\subseteq V$ such that, at time
$t=\hit/(a\ln n)$ for some universal constant $a>0$, the expected fraction of black particles in~$S$ in the
configuration~$\sigma(t)$ is at least~$\frac{1}{20}$.
\end{lemma}

\begin{proofof}{Lemma~\ref{lem:lowerbd}}
For any vertex $v\in V$, let $\tau(v)$ denote the expected hitting time in the single-particle random walk
from~$v$ to the set of heat baths. 
Thus $\hit=\max_v \tau(v)$.  For non-negative integers~$k$, define $\alpha_k:=\frac{k}{\log_c n}$, where $c>2$
is a constant to be specified later, and consider the increasing sequence of subsets $S_k\subseteq V$ defined by $$
   S_k := \left\{v\in V: \tau(v) > (1-\alpha_k)\hit\right\},\qquad 1\le k\le 1+\lfloor\log_c n\rfloor.  $$
Let $K$ be the smallest~$k$ for which $\frac{|S_{k+1}|}{|S_k|} \le c\,$; such a~$K$ must exist
since $|V|=n$.  We will take $S_{K+1}$ to be the set~$S$ whose existence is claimed in the Lemma, 
and denote by~$F$ the expected fraction of black particles in~$S$. Our goal is to show that $F\ge\frac{1}{20}$.

For $2\le k\le K+1$, define $T_k:=S_k\setminus S_{k-1}$, and $T_1:=S_1$.  Note that for any vertex $v\in T_k$ we have $$
    (1-\alpha_k)\hit < \tau(v) \le (1-\alpha_{k-1})\hit.  $$
On the other hand, for $v\in T_K$ we also have $$
    \tau(v) \le t + (1-p_0)(1-\alpha_{K+1})\hit + \sum_{r=0}^{K} q_r(1-\alpha_{K-r})\hit,   $$
where $p_0$ is the probability that the particle that was initially at~$v$, is black and lies inside $S=S_{K+1}$ at time~$t$,
and $q_r$ is the probability that the particle is black and lies in~$T_{K-r+1}$ at time~$t$.  Note that $p_0=\sum_{r=0}^K q_r $.
Combining the above two bounds yields $$
    (1-\alpha_K)\hit \le t + (1-p_0)(1-\alpha_{K+1})\hit + \sum_{r=0}^{K} q_r(1-\alpha_{K-r})\hit.  $$
Plugging in $t=\hit/(a\ln n)$ and using the fact that $p_0 = \sum_r q_r$, this becomes $$
    -\alpha_K \le \frac{1}{a\ln n}  -(1-p_0)\alpha_{K+1} - \sum_{r=0}^{K} q_r\alpha_{K-r}.  $$
Recalling that $\alpha_k=\frac{k}{\log_c n}$, we may rewrite this as $$
    -K \le \delta -(K+1)(1-p_0) - \sum_{r=0}^{K} q_r(K-r), $$
where $\delta:=(a\ln c)^{-1}$.  We assume in what follows that $K\ge 2$; if not, then a simpler version of 
the calculation below yields a better lower bound on~$F$.  Using $p_0 = \sum_r q_r$ again and rearranging gives $$
     p_0 \ge 1-\delta - \sum_{r=0}^{K} rq_r = 1 - \delta - q_1 - \sum_{r=2}^K rq_r = 1 - \delta -p_0 + q_0 - \sum_{r=2}^K (r-1)q_r,  $$
which, discarding the~$q_0$ on the RHS, implies
\begin{equation}\label{eq:pbound}
     p_0 \ge \frac{1}{2}\Bigl(1-\delta - \sum_{r=2}^K (r-1)q_r\Bigr) =\frac{1}{2}\Bigl(1-\delta - \sum_{r=2}^K p_r\Bigr),
\end{equation}     
where we have introduced the notation $p_r:=\sum_{j=r}^K q_j$ for the probability that the particle is black
and lies in~$S_{K-r+1}$ at time~$t$.  Note that this is consistent with our earlier definition of~$p_0$.
   
Consider the particles that are in $T_K$ at time $0$. Let $\hat p_r$ be the expected fraction of these particles
that at time~$t$ remain black and inside~$S_{K-r+1}$.  Note that our goal is a lower bound on~$\hat p_0$, which 
is the expected fraction of black particles remaining in~$S_{K+1}$.

Applying~\eqref{eq:pbound} to $v$ uniformly sampled in $T_K$ gives
\begin{equation}\label{eq:bbound}
    \hat p_0 \ge \frac{1}{2}\Bigl(1-\delta - \sum_{r=2}^K \hat p_r\Bigr).
\end{equation}    

Now we observe that, since each vertex contains only one particle at any time,
$\hat p_r \le \frac{|S_{K-r+1}|}{|T_K|}$.  Moreover, by our choice of~$K$,
$|S_{K-r+1}|\leq |S_K|c^{-(r-1)}$ for $r\ge 2$.  Also, $|T_K| = |S_K|-|S_{K-1}|\geq |S_K|(c-1)/c$.
Combining these observations we see that, for any $r\ge 2$,
$$
   \sum_{r=2}^K \hat p_r  \leq \frac{c}{c-1}\sum_{r=2}^K  c^{-(r-1)} 
           \leq \frac{c}{c-1} \times \frac{1}{c-1} = \frac{c}{(c-1)^2}.  $$

Plugging this bound into~\eqref{eq:bbound} yields
\begin{equation}\label{eq:bbound2}
   \hat p_0 \ge \frac{1}{2}(1-\delta -c(c-1)^{-2}).
\end{equation}   
Noting from our definition of~$K$ that 
$\frac{|T_K|}{|S|} = \frac{|S_K|-|S_{K-1}|}{|S_{K+1}|} \ge \frac{(1-1/c)|S_K|}{|S_{K+1}|} \ge \frac{c-1}{c^2}$,
and combining this bound with that on~$\hat p_0$ in~\eqref{eq:bbound2} yields the following lower bound on the desired
fraction~$F$: $$
    F\ge \hat p_0\frac{|T_K|}{|S|} \ge \frac{1}{2}(1-\delta -c(c-1)^{-2})\times \frac{c-1}{c^2}.  $$
Finally, setting $c=5$ and $\delta=\frac{1}{16}$ (which is achieved by setting the constant $a=16/\ln 5\approx 10$) 
results in $F\ge\frac{1}{20}$, as desired.
\end{proofof}

\end{document}
